지하철을 타고 등교하는 도현이는 매번 숭실대입구역에서 내린다. 숭실대입구역에서 출구로 나가기 위해서는 3개의 에스컬레이터를 타야 한다.

길이가 $L$인 에스컬레이터는 가만히 서서 올라가면 $L$초가 걸리고, 체력을 소모하지 않는다. 길이가 $L$인 에스컬레이터를 걸어서 올라가면 $\frac{L}{2}$초 만에 올라갈 수 있지만, $\frac{L}{2}$만큼의 체력을 소모한다. 단, 에스컬레이터를 타는 도중에 올라가는 방법을 바꿀 수는 없다.

도현이는 현재 $H$의 체력을 가지고 있다. 체력을 전부 소모하면 학교에 가서 공부를 할 수 없으므로 남은 체력이 $1$ 이상이 되도록 올라가려고 한다.

3개의 에스컬레이터의 길이 $A$, $B$, $C$와 체력 $H$$(1 \le A, B, C, H \le 200)$가 주어졌을 때, 도현이가 체력을 $1$ 이상 남기면서 가장 빠르게 출구로 나갔을 때 걸리는 시간을 출력하는 프로그램을 작성하라.

에스컬레이터의 길이는 항상 짝수로 주어지며, 정답은 정수임이 보장된다.